\chapter{Operating Systems in This System}
\label{app:os-topics}

This appendix relates the subtopics of final-examination topic 17 (Operating
systems) that the body does not treat to the system as built. Each item says what
the mechanism is, where it acts in \texttt{evoalgo}, and what was measured on the
recording host (SUSE Linux Enterprise Server 15 SP6, kernel 6.4, 80 logical CPUs,
755 GB of memory, no swap). Processes, threads, IPC, signals, namespaces, resource
limits, deadlock and the virtual address space are in Sections~\ref{sec:design-concurrency},
\ref{sec:design-sandbox} and~\ref{sec:impl-concurrency}.

\section{The kind of operating system the tool needs}

Linux here is a multi-user, multiprogramming, time-sharing operating system with
a monolithic kernel \citep{silberschatz2018os,tanenbaum2015modern}. The design
depends on each of these properties. Multiprogramming and time-sharing let eight
verifier workers, their runners, the engine and the UI server share the CPUs.
Being multi-user, it has per-user credentials and permission checks, which is what
makes an unprivileged user namespace safe to create: inside it a process is root
only over objects the namespace owns. The kernel, not a library, enforces the
mount, PID, network and IPC isolation; a batch or single-user system could not
host the candidate jail without root. The design also assumes two pieces of CPU
support. The first is a memory-management unit that gives each process its own
address space, so a candidate's memory cannot reach the verifier's. The second is
the user/kernel mode distinction, which makes \texttt{mount\_setattr} and
\texttt{pivot\_root} system calls the jail requests rather than operations it can
perform itself. The same kernel mechanisms confine the coding-agent harness
(Section~\ref{sec:impl-agents}).

\section{CPU scheduling}

The verifier stacks three levels of scheduling. The first two are its own; the
third is the kernel's. First, chunks are dispatched to workers first-come,
first-served in chunk order; it is not preemptive, and a hanging chunk is stopped
only by the per-chunk wall-clock timeout, which acts like a fixed time quantum
for a whole job. Second, inside the UI server, \texttt{asyncio} schedules its
tasks cooperatively: a task runs until its next \texttt{await}, so there is no
preemption and no lock is needed (Section~\ref{sec:impl-concurrency}). Third, the
kernel's fair scheduler time-shares the worker processes: each runnable process
gets CPU in proportion to its weight, the process that has run least is picked
next, and a process that waits on a pipe is not runnable and gets no share. On
this host the worker count (at most eight) is below the CPU count, so the
CPU-bound runners are rarely preempted at all.

Starvation is ruled out at each level. In-order dispatch means a chunk waits for
at most \texttt{workers} earlier ones; cooperative tasks yield at every
\texttt{await} on the socket or the file; and the fair scheduler bounds the wait
of every runnable process. Shortest-job-first would be a poor choice for the
first level even though chunk costs differ: the merge must proceed in chunk order
(Section~\ref{sec:impl-concurrency}), so finishing a short later chunk first
gains nothing until the earlier one is done.

\section{Mutual exclusion without hardware support}

Peterson's algorithm achieves mutual exclusion for two threads with two flags
and a turn variable, using only ordinary loads and stores
\citep{herlihy2012art}. It is correct only if those loads and stores are
sequentially consistent, which modern CPUs and compilers do not guarantee
without memory barriers, and it busy-waits. The system needs mutual exclusion
in two places, and Peterson's algorithm fits neither. The two sides of the
harness workspace check are separate processes with no shared memory to hold
flags, and the waiting side should sleep, not spin. The workspace therefore uses
\texttt{fcntl.flock}: the kernel holds the lock and blocks the waiter, so the
critical section (grading the snapshot and appending to \texttt{results.tsv})
is entered by one check at a time, without spinning. The append-only event and
control logs avoid a critical section altogether by making each record one
\texttt{write} (Section~\ref{sec:design-concurrency}).

\section{Paging, segmentation and the page table}

The address-space limit of the sandbox (Section~\ref{sec:design-sandbox}) acts
on paged virtual memory. With paging, a virtual address is a page number and an
offset, and the page table maps each page to a physical frame. Pages are of
fixed size (4096 bytes on this host), so physical memory has no external
fragmentation, only internal fragmentation inside a process's last page of each
mapping. Segmentation instead divides the address space into variable-sized
logical segments (code, data, stack) with a base and a limit. It matches how a
program is organised, and it protects each segment as a unit. Its cost is
external fragmentation, because segments of different sizes leave holes in
physical memory. x86-64 keeps segments only vestigially (flat segments, with
\texttt{fs} and \texttt{gs} for thread-local storage), so a runner's memory is
protected page by page: each mapping of the interpreter, NumPy and the
candidate's data is a set of pages with its own permissions.

A flat page table for a 48-bit virtual address space would need $2^{36}$
entries per process. The host's CPUs report 48-bit virtual addresses, so the
kernel uses a four-level page table: each level is indexed by 9 bits of the
address, and a lower-level table exists only for the parts of the space that
are populated. A runner that maps 4 GB of address space but touches 600 MB
therefore pays for page tables only for the ranges it populates, and, by demand
paging (Section~\ref{sec:design-sandbox}), for frames only for the pages it
touched. Every memory access would need four table look-ups. The translation
lookaside buffer (TLB) caches recent translations, so that most accesses need
none. Two effects matter here. A chunk's per-graph data is tiny (adjacency
matrices of at most five vertices, a few hundred bytes as NumPy arrays), so the
pages it touches repeatedly are few and the TLB rarely misses. A runner is also
a fresh process, which starts with no cached translations of its own; the 80
CPUs support process-context identifiers (\texttt{pcid}), which let the
translations of different processes stay in the TLB across a context switch, so
switching between the concurrently running workers and runners does not flush
it each time. Transparent huge pages are enabled on the host (\texttt{always}), so
large mappings can use 2 MB pages and one TLB entry covers 512 ordinary pages.

\section{Page replacement}

When memory is short, the kernel must choose which page to evict. Least recently
used (LRU) evicts the page unused for longest, a good approximation of the
optimal choice for programs with locality, but exact LRU needs a timestamp per
access, so kernels approximate it. Linux keeps active and inactive lists and
uses the hardware accessed bit, a clock-like second chance
\citep{silberschatz2018os}. On this host the question is mostly avoided by
design: there is no swap, so anonymous pages (the runners' heaps) cannot be
evicted at all, and under real memory pressure the kernel's out-of-memory killer
ends a process instead. That is the failure mode the verifier now handles. The
worker count is capped so that the measured demand fits the available memory
(2 GB per worker; 14.5 GB in total for eight workers at $n \le 5$), and a worker
the kernel kills anyway is detected as \texttt{BrokenProcessPool}, after which the
remaining chunks are checked serially with the same result
(Section~\ref{sec:impl-concurrency}). File-backed pages, such as the
interpreter's code and the test-graph data, can be evicted and read back, and
the read-only jail mounts share them between all runners.

\section{Input/output scheduling}

The verifier's disk I/O is small: it reads the nauty graph data once per process
and writes one generation file and a few log lines per generation. The host's
block devices use the \texttt{mq-deadline} I/O scheduler, which sorts requests to
reduce seeks and gives each request a deadline so that none waits indefinitely.
It thus prevents starvation at the device, as the fair scheduler does for the
CPU. The I/O that matters for this design is pipe and socket I/O. The sandbox's
feed and drain threads use blocking reads and writes, because each thread
serves exactly one pipe, and a blocked thread costs nothing while it waits. The
UI server serves many WebSockets on one thread, so it uses non-blocking I/O
through the \texttt{asyncio} event loop, which on Linux waits on
\texttt{epoll}: one wait covers all sockets, and a task resumes only when its
socket is ready.

\section{Storage, file systems and redundancy}

Three kinds of file system hold the system's data, chosen per use. The run
directories are on NFS (the shared project storage), a network file system whose
server owns the data and caches it on the clients. The host's \texttt{/tmp} is
btrfs, a local copy-on-write file system. The jail's \texttt{/tmp} and the
covers of the hidden directories are tmpfs, which lives in memory and disappears
with the mount namespace. The persistence design relies on two file-system
services and states where they stop holding. The first is \texttt{os.replace},
an atomic rename within one directory, which makes a generation file appear
whole or not at all. The second is \texttt{O\_APPEND} single-write appends for
the logs, which are atomic on a local file system but not guaranteed when
processes on different NFS clients append to one file. Only the engine appends
to \texttt{events.jsonl}, so there the question does not arise; the several
writers of \texttt{control.jsonl} (the CLI and the UI server) are safe as long
as they run on one host, which is how the tool is used, but the tool does not
enforce it. Every record is followed by \texttt{fsync}, so an acknowledged event
or command survives a crash of the process that wrote it. The generation file
and the \texttt{COMMITTED} marker are each replaced atomically after an
\texttt{fsync} of the file and followed by one of the directory, so the marker
never names a generation that is not on the device.

Redundancy (RAID mirroring or parity across disks, and server-side replication)
belongs to the storage layer under NFS, not to this tool. The tool's own answer
to data loss is the backup rule of Section~\ref{sec:ud-security}: a run is
exactly its directory, so a copy is a complete backup. The append-only event
log, the checksummed generation files and the commit marker let a copied or
interrupted run be verified and resumed.
